% !PRINTED 525
\weberpage{525}
Es sei jetzt $\pi_1, \pi_2, \pi_3, \ldots$ eine beliebige Anzahl von einander verschiedener gegebener Primfunctionale. Wir setzen

\[
\omega_1 = \omega\pi_2\pi_3\ldots, \quad \omega_2 = \omega\pi_1\pi_3\ldots, \quad \omega_3 = \omega\pi_1\pi_2\ldots,
\]
und bestimmen nach dem, was soeben bewiesen ist, die ganzen Zahlen $\alpha_1, \alpha_2, \alpha_3, \ldots$ in $\Omega$, so dass
\[
\begin{array}{l}
\alpha_1 \text{ theilbar wird durch } \omega_1 \text{, aber nicht durch } \omega_1\pi_1\\
\alpha_2 \quad \text{\textquotedbl} \quad \text{\textquotedbl} \quad \text{\textquotedbl} \quad \omega_2 \quad \text{\textquotedbl} \quad \text{\textquotedbl} \quad \text{\textquotedbl} \quad \omega_2\pi_2\\
\alpha_3 \quad \text{\textquotedbl} \quad \text{\textquotedbl} \quad \text{\textquotedbl} \quad \omega_3 \quad \text{\textquotedbl} \quad \text{\textquotedbl} \quad \text{\textquotedbl} \quad \omega_3\pi_3\\
\cdots\cdots\cdots\cdots\cdots\cdots\cdots\cdots\cdots\cdots\cdots\cdots\cdots\cdots\cdots\cdots,
\end{array}
\]
und leiten daraus die ganze Zahl
\[
\alpha = \alpha_1 + \alpha_2 + \alpha_3 + \cdots
\]
ab. Diese Zahl ist offenbar theilbar durch $\omega$, da alle Summanden $\alpha_1, \alpha_2, \alpha_3, \ldots$ durch $\omega$ theilbar sind. Sie ist aber nicht theilbar durch $\omega\pi_1$, weil zwar $\alpha_2, \alpha_3, \ldots$, nicht aber $\alpha_1$ durch $\omega\pi_1$ theilbar ist; und ebenso ist sie nicht durch $\omega\pi_2, \omega\pi_3, \ldots$ theilbar. Wenn wir also
\[
\alpha = \omega\eta
\]
setzen, so ist $\eta$ ein ganzes Functional, das nicht durch $\pi_1, \pi_2, \pi_3, \ldots$ theilbar ist, und das daher, wenn $\pi_1, \pi_2, \pi_3, \ldots$ die von einander verschiedenen Primfactoren von $\mu$ sind, relativ prim zu $\mu$ ist.

Nehmen wir nun beliebig eine durch $\omega$ theilbare ganze Zahl $\beta$ in $\Omega$ an und setzen $\beta = \omega\mu$, dann können wir $\alpha = \omega\eta$ so bestimmen, dass $\eta$ relativ prim zu $\mu$ wird, und dann ist $\omega$ der grösste gemeinschaftliche Theiler von $\alpha$ und $\beta$. Daraus folgt:

4. \gesperrt{Jedes ganze Functional $\omega$ des Körpers $\Omega$ ist der grösste gemeinschaftliche Theiler zweier ganzer Zahlen, und folglich ist $\omega$ associirt mit einer binären Linearform $\alpha x + \beta y$, in der $\alpha$ und $\beta$ ganze Zahlen sind.}

\begin{center}
\rule{4cm}{0.4pt}
\end{center}

Es mag hier noch eine allgemeine, auf ein anderes Gebiet hinübergreifende Bemerkung ihren Platz finden.

Es ist das Hauptergebniss dieses Abschnittes, dass sich die ganzen algebraischen Zahlen in einem bestimmten Körper in eindeutiger Weise in Primfactoren zerlegen lassen, genau in derselben Weise, wie dies bei den ganzen rationalen Zahlen bekannt ist; freilich aber nur dadurch, dass der Inhalt des Körpers
%CONT>>
